\documentclass[11pt,letterpaper]{article}
\usepackage[T1]{fontenc}
\usepackage[utf8]{inputenc}
\usepackage[scaled=0.95]{helvet}
\renewcommand{\familydefault}{\sfdefault}
\usepackage[top=1.7cm,bottom=1.5cm,left=2.2cm,right=2.2cm]{geometry}
\usepackage{microtype}
\usepackage[hidelinks]{hyperref}
\pagestyle{empty}
\setlength{\parindent}{0pt}
\setlength{\parskip}{7pt}

\begin{document}

{\large\bfseries Grover Aruquipa}\par\vspace{-6pt}
{\small Qu\'ebec, QC, Canada \ \textbullet\ \ grover.aruquipa.9@gmail.com \ \textbullet\ \ groveraruquipa.github.io}

\vspace{4pt}
Hiring Team, Apptronik, Austin\\
Mechanical Engineer, Robotic Hand Design

Dear Hiring Team,

The hardest part of a dexterous hand is not one finger. It is fitting many degrees of freedom, their actuation and their wiring into a volume the size of a human hand, while every joint survives off-axis loads and the whole hand still grasps gently. I have spent my PhD on that problem from the mechanism side, and I would like to bring it to the hand of Apollo.

I am finishing my PhD at Universit\'e Laval with Prof.\ Cl\'ement Gosselin, whose lab has designed robotic hands and grippers for more than twenty years, and Prof.\ Patrice Lambert. I defend in early 2027. I designed, built and tested three high-DoF robots with 4, 7 and 8 actuators in which the gripper is driven remotely, from motors at the base, through linkages, concentric hollow shafts and belts. For each one I defined the architecture and the load paths, did the SolidWorks design and the FEA, selected the motors, bearings and encoders, built the prototype, integrated it with the electronics and the control, and tested it until it worked.

Packaging is where I have spent most of my design hours. My 4-DoF coaxial wrist-gripper fits four actuation channels into concentric hollow shafts in a 180\,mm footprint with 0.25\,kg of moving mass, and it turned its tool ten full revolutions with less than one degree of tracking error while holding the object. I am now building a Delta robot with an integrated Sarrus-linkage gripper that estimates the grasp force from the motor currents, with no sensor in the fingers, to pick fragile objects quickly. I also wrote a published review in the ASME Journal of Mechanisms and Robotics on parallel mechanisms for grasping, including parallel and underactuated robot hands, so I know the trade-offs between finger topologies, remote and local actuation, and underactuation.

I work closely with controls and software because on my own robots I was also the one writing the 500\,Hz real-time control and the kinematics, and I lead two graduate students day to day as technical lead of an 8-actuator robot. I would be glad to bring that ownership to the hand team at whatever level fits best.

I am available from January 2027 and ready to relocate to Austin. Prof.\ Gosselin, Prof.\ Lambert and Univ.-Prof.\ Andreas M\"uller of JKU Linz can speak to my work; their details are in my CV.

Thank you for your time. I would be happy to walk you through the prototypes on video or in person.

Sincerely,\\
Grover Aruquipa

\end{document}
